\section{Tabellen}

\begin{table}[ht]
	\centering
	\caption{Messwerte der Probe A}
	\begin{tabular}{l S[table-format=3.1] S[table-format=2.2]}
		\toprule
		Probe & {Länge (\si{\cm})} & {Masse (\si{\gram})} \\
		\midrule
		1     & 12.3               & 4.56                 \\
		2     & 11.8               & 4.12                 \\
		3     & 12.5               & 4.78                 \\
		\bottomrule
	\end{tabular}
\end{table}

\begin{table}[ht]
	\centering
	\caption{Messwerte der Probe A}
	\csvautotabular{data/messwerte.csv}
\end{table}

\begin{table}[ht]
	\centering
	\caption{Messwerte formatiert mit siunitx}
	% S[table-format=...] tells siunitx how to align numbers (decimal points)
	\begin{tabular}{l S[table-format=2.1] S[table-format=1.2]}
		\toprule
		Probe    & {Länge (\si{\cm})} & {Masse (\si{\gram})} \\
		\midrule
		\csvreader[
		late after line=                                     \\,
			head to column names
		]{data/messwerte.csv}{}%
		{%
		\csvcoli & \num{\csvcolii}    & \num{\csvcoliii}     %
		}
		\bottomrule
	\end{tabular}
\end{table}
